\documentclass[10pt]{article}

\usepackage[margin=0.7in]{geometry}
\usepackage{amsmath}
\usepackage{amssymb}
\usepackage{graphicx}
\usepackage{booktabs}
\usepackage[font=small,skip=2pt]{caption}

\setlength{\parindent}{0pt}
\setlength{\parskip}{4pt}
\pagestyle{empty}

\begin{document}

\subsection*{Lab 2 (a)--(b): inner-loop controller, revised for the measured motor time constant}

\textbf{Model update.} Part (c) captures with the original $C_1$ (10 steps,
$-0.7 \leftrightarrow 0.7$ rad) overshoot by $5.8\%$ against $\approx 1\%$
predicted. The rise matches the Lab 1 model, so $K_1$ holds; the extra coasting
near the target indicates a slower motor. A least-squares fit of the sampled
loop (Tustin $D_1$, $0.3$ V stiction offset, Coulomb friction $F$) to the mean
response over $(K_1, \tau, F)$ gives $K_1 = -1.81$ (unchanged),
$\boldsymbol{\tau = 0.0231}$ \textbf{s} ($1.4\times$ Lab 1) and $F \approx 0.05$ V,
with $11\times$ lower squared error than $\tau = 0.0165$ s (Fig.~1a). The old
controller zero ($-64.4$) then misses the motor pole ($-1/\tau = -43.3$), and
even the linear model no longer meets (a): $t_s = 0.37$ s (Table~1). Parts
(a)--(b) are redone with $\tau = 0.0231$ s.

\textbf{(a) Pole placement.} $G_p = \frac{K_1}{s(\tau s + 1)}$,
$C_1 = \frac{b_1 s + b_0}{s^2 + a_1 s + a_0}$, closed-loop poles
$-\zeta\omega_n \pm j\omega_d$, $-p_3$, $-1/\tau$ with $\zeta = 0.8$,
$\zeta\omega_n = 15.5$, $p_3 = 2\zeta\omega_n = 31$. The pole at $-1/\tau$ forces
$b_1 = \tau b_0$ (controller zero cancels the motor pole), leaving
\begin{gather*}
s^3 + a_1 s^2 + a_0 s + K_1 b_0 = (s^2 + 2\zeta\omega_n s + \omega_n^2)(s + p_3) \\
a_1 = 2\zeta\omega_n + p_3 = 62,\quad a_0 = \omega_n^2 + 2\zeta\omega_n p_3 = 1336.390625,\quad
b_0 = \frac{\omega_n^2 p_3}{K_1} = -6429.342196,\quad b_1 = \tau b_0
\end{gather*}
\[
\boxed{C_1(s) = \frac{-148.5178047\, s - 6429.342196}{s^2 + 62\, s + 1336.390625}}
\qquad
T(s) = \frac{11637.11}{s^3 + 62 s^2 + 1336.39 s + 11637.11},\;\; T(0) = 1 .
\]
Only $b_1$ depends on $\tau$, so the closed-loop poles and $\theta(t)$ are as
before; the slower motor raises the peak voltage from $5.66$ to $5.85$ V.

\textbf{(b) Discretization}, $T_s = 1$ ms ($|s| T_s \le 0.044$ for all controller roots):
\begin{align*}
D_{1,\text{Tustin}}[z] &= \frac{-0.07356126214\, z^2 - 0.003117002633\, z + 0.07044425951}{z^2 - 1.938587902\, z + 0.9398836903} \\
D_{1,\text{ZOH}}[z] &= \frac{-0.1471242781\, z + 0.1408908837}{z^2 - 1.938587225\, z + 0.9398828868}
\end{align*}
Implemented (Tustin, Appendix 2 form):
\[
u = 1.938587902049457\,u_1 - 0.9398836903183\,u_2 - 0.073561262142669\,e - 0.003117002633164\,e_1 + 0.070444259509506\,e_2
\]

\vspace{-4pt}
\begin{figure}[h]
\centering
\includegraphics[width=\textwidth]{lab2_revision.pdf}
\caption{(a) Part (c) hardware steps with the original $C_1$ (band: all steps)
against the sampled-loop model with $\tau = 0.0165$ and $0.0231$ s; dotted line:
$5\%$ OS. (b), (c) Revised design on the $\tau = 0.0231$ s plant.}
\end{figure}

\vspace{-10pt}
\begin{table}[h]
\centering
\small
\caption{Step $-0.7 \to 0.7$ rad. Specs: (a)--(b) $OS \le 5\%$, $t_s \le 0.25$ s,
$|V| \le 6$ V; (c) $t_s \le 0.40$ s.}
\begin{tabular}{llccc}
\toprule
Controller & Plant / implementation & OS & $t_s$ (2\%) & Peak $|V|$ \\
\midrule
Original $C_1 = \frac{-102 s - 6564}{s^2 + 61.4 s + 1359}$ & $\tau = 0.0165$, continuous & $0.96\%$ & $0.241$ s & $5.66$ V \\
 & $\tau = 0.0231$, continuous & $3.24\%$ & $0.373$ s & $5.77$ V \\
 & $\tau = 0.0231$, Tustin & $3.42\%$ & $0.376$ s & $5.78$ V \\
 & hardware, $0.3$ V stiction (runs 2--3) & $5.3$--$6.1\%$ & $0.34$--$0.45$ s & $5.98$--$6.03$ V \\
\midrule
Revised $C_1$ & $\tau = 0.0231$, continuous & $1.02\%$ & $0.243$ s & $5.85$ V \\
 & $\tau = 0.0231$, Tustin & $1.15\%$ & $0.240$ s & $5.86$ V \\
 & $\tau = 0.0231$, ZOH & $1.28\%$ & $0.239$ s & $5.88$ V \\
 & HW model, Tustin, $0.3$ / $0.15$ V stiction & $3.3$ / $1.9\%$ & $0.31$ / $0.23$ s & -- \\
\bottomrule
\end{tabular}
\end{table}

\end{document}
